\subsection{平展且分离映射的纤维基数的上界}

定理 3. --- 设 E 与 B 为拓扑空间，p: E → B 为一个平展且分离的映射。假设 E 是连通的，且 B 是局部连通的。设 W 为 B 的拓扑的一个基。那么，对 B 的每一点 a，有 Card(E \(_{a}\) ) ≤ sup(Card(W), Card(N))。

由于空间 B 是局部连通的，它的拓扑有一个由连通开集组成的基 \(\mathcal{W}'\) ，例如诸开集 \(\mathcal{W}\) 的连通分支 (TG, I, p. 85, prop. 11)。根据下文的引理 4，存在由连通开集组成的 B 的拓扑的一个基 \(\mathcal{V}\) ，且满足 \(\mathrm{Card}(\mathcal{V}) \leqslant \mathrm{Card}(\mathcal{W})^2\) 。

引理 4. --- 设 B 为一个拓扑空间，并设 W 与 W' 为 B 的拓扑的基。存在 W' 的一个子集 V，它是 B 的拓扑的一个基，且满足 \(\text{Card}(\mathcal{V}) \leqslant \text{Card}(\mathcal{W})^{2}\) 。

设 \(\mathcal{A} \subset \mathcal{W} \times \mathcal{W}\) 为如下的对 \((\mathrm{W}_1, \mathrm{W}_2)\) 所成的集：存在 \(\mathrm{W}' \in \mathcal{W}'\) 使 \(\mathrm{W}_1 \subset \mathrm{W}' \subset \mathrm{W}_2\) ；设 \(\varphi: \mathcal{A} \to \mathcal{W}'\) 为一个映射，使得 \(\mathrm{W}_1 \subset \varphi(\mathrm{W}_1, \mathrm{W}_2) \subset \mathrm{W}_2\) 对每个对 \((\mathrm{W}_1, \mathrm{W}_2) \in \mathcal{A}\) 成立，并设 \(\mathcal{V} \subset \mathcal{W}'\) 为 \(\varphi\) 的像。我们有

\[
\operatorname{Card} (\mathcal {V}) \leqslant \operatorname{Card} (\mathcal {A}) \leqslant \operatorname{Card} (\mathcal {W} \times \mathcal {W})
\]

(E, III, p. 25, prop. 3)。我们来证明 V 是 B 的拓扑的一个基。设 x 为 B 的一点，U 为 x 的一个邻域。由假设，x 有一个邻域 \(W_{2} \in W\) 含于 U，有一个邻域 \(W' \in W'\) 含于 \(W_{2}\) ，且有一个邻域 \(W_{1} \in W\) 含于 \(W'\) 。于是有 \(W_{1} \subset W' \subset W_{2} \subset U\) 。那么，\((W_{1}, W_{2}) \in \mathscr{A}\) ，且 \(\varphi(W_{1}, W_{2})\) 是 x 的一个含于 U 的邻域。根据 TG, I, p. 5 的 prop. 3，集 V 是 B 的拓扑的一个基。

设 V 为由非空连通开集组成的 B 的拓扑的一个基，且满足 \(\text{Card}(\mathcal{V}) \leqslant \text{Card}(\mathcal{W})^{2}\) 。设 U 为 E 的如下的开集 U 所成的集：p 诱导从 U 到属于 V 的某个开集 V 上的一个同胚。U 的每个元素都是非空连通开集。由于映射 p 是平展的，根据定义 2 (I, p. 28) 与 TG, I, p. 5 的 prop. 3，集 U 是 E 的拓扑的一个基。我们把 U 的元素的每个有限序列 \((\mathrm{U}_{0}, \ldots, \mathrm{U}_{n})\) 称为毛虫，如果对 \(1 \leqslant i \leqslant n\) ，\(U_{i-1} \cap U_{i}\) 非空且 \(p(\mathrm{U}_{i-1}) \cap p(\mathrm{U}_{i})\) 连通。用 S 表示 V 的元素的有限序列所成的集，并对每个毛虫 \(c = (\mathrm{U}_{0}, \ldots, \mathrm{U}_{n})\) ，用 \(p(c)\) 表示序列 \((p(\mathrm{U}_{0}), \ldots, p(\mathrm{U}_{n}))\) 。

引理 5. --- a) 设 \(c = (\mathrm{U}_0, \ldots, \mathrm{U}_n)\) 与 \(c' = (\mathrm{U}_0', \ldots, \mathrm{U}_m')\) 为两个毛虫，满足 \(\mathrm{U}_0 = \mathrm{U}_0'\) 且 \(p(c) = p(c')\) 。那么 \(c = c'\) 。

\begin{enumerate}
\def\labelenumi{\alph{enumi})}
\setcounter{enumi}{1}
\item
  设 U 与 U' 为 U 的元素。则存在一个毛虫 \(c = (\mathrm{U}_{0}, \ldots, \mathrm{U}_{n})\) ，使得 \(U_{0} = U\) 且 \(U_{n} = U'\) 。
\item
  必然有 m = n。对每个满足 \(0 \leqslant i \leqslant n\) 的整数 i，令 \(V_{i} = p(U_{i})\) ，并用 \(s_{i}\) 与 \(s_{i}^{\prime}\) 表示 p 在 \(V_{i}\) 上的、其像分别为 \(U_{i}\) 与 \(U_{i}^{\prime}\) 的连续截面。为证等式 \(c = c^{\prime}\) ，我们对 i 用归纳法证明等式 \(s_{i} = s_{i}^{\prime}\) （对每个满足 \(0 \leqslant i \leqslant n\) 的整数 i 成立）。由假设，有 \(s_{0} = s_{0}^{\prime}\) 。设 i 满足 \(1 \leqslant i \leqslant n\) 且 \(s_{i-1} = s_{i-1}^{\prime}\) 。由于 \(U_{i-1} \cap U_{i}\) 含有一点 x，连续截面 \(s_{i-1}\) 与 \(s_{i}\) 在 \(p(x)\) 处重合，从而在 \(V_{i-1} \cap V_{i}\) 的每一点处重合，因为它是连通的 (I, p. 34 的推论 1)。对 \(s_{i-1}^{\prime}\) 与 \(s_{i}^{\prime}\) 同理。出于同样的理由，\(s_{i}\) 与 \(s_{i}^{\prime}\) （在 \(V_{i-1} \cap V_{i}\) 中重合者）也在 \(V_{i}\) 中重合，由此即得结论。
\item
  若 \(x\) 与 \(y\) 为 E 的点，则称毛虫 \(c = (\mathrm{U}_0, \ldots, \mathrm{U}_n)\) 连接 \(x\) 到 \(y\) ，如果 \(x \in \mathrm{U}_0\) 且 \(y \in \mathrm{U}_n\) 。在集合 E 中，设 R 为关系“存在连接 \(x\) 到 \(y\) 的毛虫”。关系 R 是自反的，因为 \(\mathcal{U}\) 是 E 的一个覆盖。它显然是对称的。我们来证明它是传递的：设 \(x, y, z\) 为 E 的三点，\((\mathrm{U}_0, \ldots, \mathrm{U}_n)\) 与 \((\mathrm{U}_0', \ldots, \mathrm{U}_m')\) 为分别连接 \(x\) 到 \(y\) 与 \(y\) 到 \(z\) 的毛虫。设 \(\mathrm{U} \in \mathcal{U}\) 为 \(y\) 的一个含于 \(\mathrm{U}_n \cap \mathrm{U}_0'\) 的邻域；我们有 \(p(\mathrm{U}_n) \cap p(\mathrm{U}) = p(\mathrm{U}_0') \cap p(\mathrm{U}) = p(\mathrm{U})\) ，且由于 U 连通，序列 \((\mathrm{U}_0, \ldots, \mathrm{U}_n, \mathrm{U}, \mathrm{U}_0', \ldots, \mathrm{U}_m')\) 是连接 \(x\) 到 \(z\) 的毛虫。R 的等价类都是开的，从而也是闭的。由于 E 是连通的，由此得出 E 的任意两点总可由一条毛虫连接。
\end{enumerate}

设 U 与 U′ 为 U 的元素，x 为 U 的一点，\(x'\) 为 U' 的一点。存在一个毛虫 \((\mathrm{U}_{2},\ldots,\mathrm{U}_{n-2})\) ，它连接 x 到 \(x'\) 。设 \(U_{1}\) 与 \(U_{n-1}\) 为 U 的开子集，使得 \(x \in U_{1}\) 、\(x' \in U_{n-1}\) 、\(U_{1} \subset U \cap U_{2}\) 、\(U_{n-1} \subset U' \cap U_{n-2}\) 。令 \(U_{0} = U\) ，\(U_{n} = U'\) 。那么序列 \((\mathrm{U}_{0},\ldots,\mathrm{U}_{n})\) 是一个毛虫。

现在来证明定理。设 U 为 \(\mathcal{U}\) 的一个元素，并设 C 为毛虫 \((\mathrm{U}_0,\dots ,\mathrm{U}_n)\) （满足 \(\mathrm{U}_0 = \mathrm{U}\) 者）所成的集。把 C 中的毛虫 \(c = (\mathrm{U}_0,\dots ,\mathrm{U}_n)\) 映为序列 \(p(c) = (p(\mathrm{U}_{0}), \ldots, p(\mathrm{U}_{n}))\) 的、从 C 到 S 的映射是单射（引理 5, (a)），从而有 (1) \(\text{Card}(\mathrm{C}) \leqslant \text{Card}(\mathrm{S}).\)

从 C 到 \(\mathcal{U}\) 的、把 \(c = (\mathrm{U}_0,\dots ,\mathrm{U}_n)\in \mathrm{C}\) 映为开集 \(\mathrm{U}_n\) 的映射是满射（由引理 5 的第二部分），从而

\[
\operatorname{Card} (\mathcal {U}) \leqslant \operatorname{Card} (\mathrm{C}).\tag{2}
\]

对 E 的每一点 \(x\) ，选取开集 \(\mathrm{U}_x\) （\(\mathcal{U}\) 中的），使得 \(x \in \mathrm{U}_x\) 。若 \(x\) 与 \(y\) 是同一纤维 \(\mathrm{E}_a\) （\(p\) 的）中的不同点，则有 \(\mathrm{U}_x \neq \mathrm{U}_y\) ，因为 \(p \mid \mathrm{U}_x\) 是单射。于是，对 \(a \in \mathrm{B}\) ，我们有

\[
\operatorname{Card} \left(\mathrm{E} _ {a}\right) \leqslant \operatorname{Card} (\mathcal {U}).\tag{3}
\]

最后，若 V 是有限集，则 V 的元素的有限序列所成的集 S 是可数的 (E, III, p. 49, prop. 1)，从而

\[
\operatorname{Card} (\mathrm{S}) \leqslant \operatorname{Card} (\mathbf {N}).\tag{4}
\]

若 \(\mathcal{V}\) 是无限的，我们有

\[
\operatorname{Card} (\mathrm{S}) = \operatorname{Card} (\mathcal {V}) \leqslant \operatorname{Card} (\mathcal {W}) ^ {2} = \operatorname{Card} (\mathcal {W})\tag{5}
\]

由 E, III, p. 50 的推论与 E, III, p. 49 的推论 1 即得。定理由上述关系式 (1) 至 (5) 得出。

注. --- 由前所证，若 B 的拓扑有可数基，则 E 的拓扑亦有可数基，且 E 的纤维是可数的（参看 TG, I, p. 88, Poincaré-Volterra 定理）。

